\chapter{The complete render loop, annotated}
\label{ch:loopannotated}

\begin{aside}
Every line of the main loop, with commentary. The whole thing
is maybe 80 lines of code; understanding it is understanding
\maya{}.
\end{aside}

\section{The pseudocode}

\begin{lstlisting}
void Runtime::run() {
    // [1] startup bundle
    backend.write(boot_ansi());

    // [2] initial render
    schedule_frame();

    // [3] the loop
    while (running) {
        auto deadline = next_deadline();
        auto event = backend.poll(deadline);

        // [4] dispatch and drain
        if (event) dispatch(*event);
        drain_input();

        // [5] microtasks
        run_microtasks();

        // [6] effects
        run_dirty_effects();

        // [7] render if needed
        if (frame_dirty) render_frame();
    }

    // [8] shutdown bundle
    backend.write(shutdown_ansi());
}
\end{lstlisting}

\section{[1] Startup}

Boot bundle = alt-screen + raw mode + mouse + CSI-u +
clear + home. See Chapter~\ref{ch:altscreen}.

\section{[2] Initial render}

Mark the first frame dirty so something paints
immediately.

\section{[3] Loop}

\code{deadline} is the next timer or animation tick. The
poll sleeps until input or deadline.

\section{[4] Dispatch}

Parse the event; route to the focused element. Drain any
queued events (input came faster than wake).

\section{[5] Microtasks}

Deferred functions scheduled by handlers.

\section{[6] Effects}

Dirty signals fire their effects in topological order.

\section{[7] Render}

If any signal write or input dispatch dirtied the frame,
render: layout, paint, diff, emit.

\section{[8] Shutdown}

Reverse of boot. Atexit hook backs this up.

\section{Recap}

\begin{recap}
\begin{itemize}[leftmargin=*]
  \item One loop, eight steps.
  \item Sleep until something wants the CPU.
  \item Render only when dirty.
\end{itemize}
\end{recap}

\section*{Workshop}

\begin{enumerate}[leftmargin=*]
  \item Trace your app through the eight steps for one
    keystroke.
  \item Identify where time goes.
  \item Add a custom step (e.g., a metrics flush).
\end{enumerate}
